\chapter{Iceberg Partitioning and Layout}
\label{ch:iceberg-layout}

\begin{chapterintro}
Iceberg partitioning is hidden: you filter on \texttt{event\_ts} and Iceberg
derives the predicate on \texttt{day(event\_ts)} for you, so a query never
references a partition column directly and pruning does not break when the
scheme changes. This chapter covers the partition transforms, evolving a spec
without rewriting history, sort orders and Z-ordering, target file size, and the
write-distribution mode --- the choices that decide how much a reader and a
writer have to move.
\end{chapterintro}

\section{Hidden partitioning and partition transforms}
\label{sec:ice-partitioning}

A classic Hive table's partition column is a real column duplicated into the
directory path --- a query has to filter on that exact derived value
(\texttt{WHERE event\_date = '2026-09-20'}) for pruning to fire, and a filter
on the underlying timestamp (\texttt{WHERE event\_ts >= \ldots}) does not
automatically resolve to it. Iceberg's partition spec instead names a
\emph{transform} of a source column --- \texttt{day(event\_ts)},
\texttt{bucket(16, route\_id)}, \texttt{truncate(4, zone)} --- and stores the
computed value in the manifest, not as a queryable column at all. The planner
derives the partition predicate from a filter on the source column itself, so
the specimen query's plain \texttt{event\_ts >= TIMESTAMP '\ldots'} filter
from \Cref{sec:pruning-static} prunes correctly with no partition column ever
appearing in the query text.

\Cref{tab:ice-transforms} covers the transforms this stack uses.

\begin{table}[htbp]
\centering
\caption{Common Iceberg partition transforms}
\label{tab:ice-transforms}
\begin{tabular}{@{}llp{5.2cm}@{}}
\toprule
\textbf{Transform} & \textbf{Use case} & \textbf{Notes} \\
\midrule
\texttt{day} / \texttt{hour} & Time-series data filtered by a timestamp range &
  Coarser (\texttt{day}) means fewer, larger partitions; \texttt{hour} suits
  high-volume streams (\Cref{ch:streaming}) \\
\texttt{bucket(n, col)} & High-cardinality key with equality or join
  predicates & Even file distribution by hash; enables storage-partitioned
  joins (\Cref{sec:joins-spj}) when both sides share the spec \\
\texttt{truncate(n, col)} & String or numeric prefix filters & Groups by a
  value prefix without one partition per distinct value \\
\texttt{identity} & Low-cardinality column filtered directly & Same as
  classic partitioning, just recorded the hidden way \\
\bottomrule
\end{tabular}
\end{table}

\begin{tip}
Prefer \texttt{day} over \texttt{identity} on a raw timestamp: an
\texttt{identity} partition on \texttt{event\_ts} would create one partition
\emph{per distinct timestamp value}, which is not partitioning at all --- it is
one file per row's worth of metadata. Always partition on a transform coarse
enough to group many rows per partition.
\end{tip}

\section{Partition evolution}
\label{sec:ice-partition-evolution}

A classic Hive table's partitioning is fixed at creation; changing it means
rewriting every file under the new scheme. Iceberg's partition spec can
change (\texttt{ALTER TABLE \ldots ADD PARTITION FIELD}) without touching a
single existing file, because the transform that produced a file's partition
value is recorded per-manifest, not assumed globally --- old manifests keep
describing old files under the old spec, new manifests describe new files
under the new spec, and a query spans both by evaluating each manifest
against its own recorded spec.

This is why moving from \texttt{month(event\_ts)} to \texttt{day(event\_ts)}
as a table grows is a metadata-only operation, not a maintenance project: new
data lands under the finer scheme immediately, old data keeps the coarser
scheme it was written with, and a query filtering by date prunes correctly
against both --- coarser pruning on the old partitions, finer pruning on the
new ones, with no explicit migration step and no downtime.

\section{Sort orders and Z-ordering}
\label{sec:ice-sort-order}

A partition spec controls which \emph{files} a filter can skip; a
\emph{sort order} controls how much of the \emph{content} within a kept file
can still be skipped, via the column min/max statistics
\Cref{ch:iceberg-spec}'s manifests already carry per file. Files written
sorted by a column have tight, non-overlapping min/max ranges across files ---
a filter on that column can then skip whole files even within one partition,
the same statistics-based pruning \Cref{sec:c-pushdown} described for row
groups, just at the file level.

A single sort order only helps predicates on the columns it is sorted by,
front to back. \emph{Z-ordering} (\texttt{rewrite\_data\_files} with
\texttt{sort\_order} type \texttt{zorder}) interleaves the bits of several
columns into one sort key instead, trading perfect single-column ordering for
imperfect-but-useful ordering on every column in the set simultaneously ---
the right choice when queries filter on different columns unpredictably
rather than one column consistently.

\section{Target file size and the small-files problem}
\label{sec:ice-file-size}

Every file a reader opens costs a fixed overhead --- opening the file, reading
its footer, resolving its schema --- independent of how many rows it holds
(\Cref{sec:shape-jobs}'s input-partition sizing runs into the same fixed cost
from the other direction). A table accumulating many small files, typically
from frequent low-volume commits (a streaming job flushing every micro-batch,
\Cref{ch:streaming}), pays that overhead over and over for little data each
time. \texttt{write.target-file-size-bytes} (512~MB by default) sets what a
write aims for; \Cref{ch:iceberg-write}'s compaction procedures are the
after-the-fact fix for files that already landed too small.

\section{Write distribution modes}
\label{sec:ice-write-distribution}

Before a write's rows are grouped into per-partition files, Spark may need to
shuffle them so that all rows for a given partition value land together
rather than scattered across every task's output --- exactly the shuffle
mechanics from \Cref{ch:shuffle}, applied to a write instead of a read.
\texttt{write.distribution-mode} chooses how:

\begin{description}
  \item[\texttt{none}] writes rows in whatever order the upstream plan already
        produced them, with no shuffle. Cheapest, but every task can touch
        every partition, so a table with many partitions ends up with many
        small files per commit --- one per (task, partition) pair touched.
  \item[\texttt{hash}] (the default for unpartitioned or partitioned tables
        with a defined spec) shuffles by partition value first, so each task
        after the shuffle owns one partition's rows and writes one file for
        it. This is the direct fix for the \texttt{none} mode's small-files
        problem, at the cost of the shuffle itself.
  \item[\texttt{range}] shuffles by a sampled range of the sort order instead
        of a hash, producing files with tighter, non-overlapping value
        ranges --- the write-time complement to \Cref{sec:ice-sort-order}'s
        pruning benefit, at a higher planning cost (a sampling pass before
        the real shuffle).
\end{description}

\section{Summary}
\label{sec:ice-layout-summary}

Iceberg partitioning is hidden behind a transform of a source column, so
pruning works from a plain filter on that column with no partition column in
the query; \texttt{day}/\texttt{hour} suit time-series filters,
\texttt{bucket} suits high-cardinality keys and storage-partitioned joins,
and an \texttt{identity} transform on a high-cardinality column is a
modeling mistake, not a partitioning choice. Evolving the spec is
metadata-only because old files keep the spec they were written under. A
sort order or Z-order tightens the per-file statistics that let a reader
skip whole files inside a kept partition; target file size and the write
distribution mode are what decide whether a write produces a few
well-sized files or many small ones in the first place.

\exercisehead
\begin{exercises}
\item Given a daily data volume and a query pattern, choose a partition spec.
      Justify it against the target of 128~MB--1~GB per partition per write.
\item Evolve a table's partition spec from \texttt{month(event\_ts)} to
      \texttt{day(event\_ts)}. Confirm from the \texttt{partitions} metadata
      table that old and new data report different specs.
\item Write the same DataFrame with \texttt{write.distribution-mode} set to
      \texttt{none} and then \texttt{hash}. Compare the file count per
      partition in each case.
\end{exercises}
